\section{Related Work}
\label{sec:related-work}

\subsection{SPC reconstruction, observing geometry, and morphological diversity}

SPC developed as a joint tool for surface characterization and spacecraft navigation. \citet{gaskell2008} described the use of local topography and albedo maps to construct global and regional shape products, with applications to several small bodies. \citet{gaskell2023} subsequently documented the mathematical foundations of the OSIRIS-REx implementation, including the estimation of surface and observing parameters from stereo and brightness information. The practical treatment by \citet{palmer2022} explains how suitable images, local terrain products, and successive refinements are organized into a mission workflow. These studies establish that reconstruction quality is closely connected to the geometry, resolution, and consistency of the available observations.

Morphological diversity is already present in successful SPC applications. The reconstruction of comet 67P from Rosetta imaging, for example, demonstrates applicability to a strongly non-spherical target \citep{jorda2016}. Accordingly, a claim of generalization must address more than the ability to process an irregular mesh. The relevant question for an automatic system is whether its initialization, quality assessment, and observation selection remain effective as visibility, illumination, and local terrain structure change. An observing strategy that works on an exposed convex region may leave a concavity poorly constrained, even when both regions have similar nominal sampling. This motivates evaluating generalization at the level of regional reconstruction quality and the amount of target-specific intervention required.

Dedicated SPC sensitivity tests give a stronger basis for image selection than geometric coverage alone. \citet{palmer2024geometry} studied observation combinations and reconstruction iterations, finding that a small set of complementary topographic and albedo images can produce high-quality terrain under their test conditions. The companion accuracy study examined the effects of image and processing conditions and identified weaknesses near steep terrain \citep{palmer2024accuracy}. These findings supply physical and empirical guidance for a planner. Their translation into a campaign for an unfamiliar body still requires deciding which regions need which observations, whether those observations are reachable, and how priorities should change after the next model update.

\subsection{From reconstruction diagnostics to verifiable accuracy}

The interpretation of reconstruction quality has been studied directly in the SPC literature. \citet{weirich2022} reconstructed a synthetic asteroid and compared it with known reference topography, testing the relationship between internal consistency and absolute geometric accuracy. Their results support the use of diagnostics when their interpretation has been established for the relevant conditions. They do not provide a universal conversion from a consistency statistic to a geometric error bound for every body, illumination regime, or reconstruction configuration. A planner that treats a quality score as an accuracy estimate therefore inherits a calibration problem.

The Bennu validation work of \citet{alasad2021} is particularly relevant to repeated planning rounds. It assessed successive SPC models using internal diagnostics, image comparisons, and independent laser-altimeter information. The study shows how operational model updates can be accompanied by documented quality assessments. It also provides a precedent for combining evidence with different sensitivities to error. Within an adaptive planner, these assessments would need to become spatially explicit decisions: which region requires another viewing direction, which discrepancy warrants further checking, and which area already satisfies the mission requirement?

Several distinctions are important when making that transition. Image sampling and model resolution describe the scale at which information is represented; geometric accuracy describes agreement with the actual surface. Internal consistency measures agreement among estimates or observations; it may not expose a bias shared by them. An image-based validation test also has its own dependencies on registration and reflectance modeling. These considerations motivate validating per-round indicators against known geometry in simulation and, where available, independent measurements. Held-out observations can provide additional consistency evidence, but should not be interpreted as direct ground truth. The role of the accuracy indicator in the present problem is both retrospective and prospective: it must describe the current model and support a defensible prediction of what another observation could improve.

\subsection{Point-set learning and continuous camera-extrinsic prediction}

A surface representation used for planning must accommodate irregular geometry and changes in the reconstructed model. Point clouds offer a direct description of sampled surface locations, with associated attributes when available. PointNet introduced shared pointwise processing and symmetric aggregation for learning from unordered point sets, with classification and segmentation as its original applications \citep{qi2017pointnet}. Its treatment of input order is attractive for a planning system whose sampled surface can change between reconstruction rounds. PointNet++ subsequently introduced hierarchical learning of local geometric structure and mechanisms for handling nonuniform sampling density \citep{qi2017pointnetpp}. This development highlights the relevance of both global shape and local detail when interpreting complex surfaces.

These representation properties motivate our use of a modified PointNet for predicting continuous camera extrinsics. A point-based network can relate the configuration of the currently reconstructed surface to an observation pose without depending on a common mesh connectivity across all targets. The prediction task concerns the camera's desired position and orientation. It consequently differs from assigning an object category, labeling each point, or predicting a scalar score for every member of a fixed viewpoint catalogue. This continuous pose-prediction role defines the adaptation considered here. The representation literature further motivates retaining geometric information at the spatial scales relevant to the reconstruction task.

The relation to learned NBV prediction is especially useful. PC-NBV learns candidate-view gains from a partial point cloud and the current selection state \citep{zeng2020}. Our model instead produces continuous parameters describing the next observation geometry. This opens the possibility of adjusting a pose between previously enumerated viewpoints, including changes that improve a baseline or reveal a partly occluded region. The representation must nevertheless preserve the physical information needed to interpret that pose. Normalizing the shape may assist learning, but physical scale, orientation relative to the Sun, and camera range remain necessary for evaluating resolution, illumination, and operational feasibility. Invariance to the order of points should therefore be distinguished from invariance to physically meaningful changes of observing geometry.

Point-based processing also leaves substantive generalization questions. Sparse sampling may omit a narrow neck or a small depression, and a coarse SPC model may place an occluding boundary incorrectly. A network can be insensitive to point ordering while remaining sensitive to these errors. Evaluation should consequently distinguish transfer to new bodies from robustness to incomplete or distorted reconstructions of the same body. In the proposed framework, geometric learning supports adaptation of the observation poses; the reliability of the resulting SPC surface must be checked independently through the reconstruction and validation loop.

\subsection{Autonomous estimation and active three-dimensional reconstruction}

Autonomous small-body reconstruction and navigation have advanced through joint estimation. \citet{baldini2018} introduced an optimization-based SPC-SLAM approach for estimating an unknown body's shape and motion together with spacecraft-relative navigation, using simulated observations of comet 67P. More recently, Photoclinometry-from-Motion (PhoMo) combined keypoint-based structure from motion with photometric estimation in a factor-graph formulation, jointly considering spacecraft pose, landmark position, Sun-relative direction, surface normal, and albedo \citep{driver2025}. These methods strengthen the estimation side of autonomy. They also make clear that surface quality depends on several coupled unknowns, so the information value of an image cannot generally be inferred from visibility alone. The present planning problem concerns how subsequent acquisitions should be chosen using the state of such a reconstruction process.

In robotics, next-best-view (NBV) methods provide established mechanisms for selecting informative measurements. PC-NBV predicts candidate-view information gains directly from a partial point cloud and the current view-selection state \citep{zeng2020}. SCONE uses occupancy and visibility prediction with volumetric integration to estimate gains in surface coverage, enabling reasoning about unobserved geometry \citep{guedon2022}. GenNBV studies policy generalization through reinforcement learning and geometric, semantic, and action representations \citep{chen2024}. These methods offer useful approaches to state representation and candidate evaluation. Their reconstruction objectives and sensing assumptions, however, do not establish how their scores translate into SPC terrain accuracy under varying solar illumination.

The choice between a single next view and a group of views has also received attention. AIR-OSVP combines implicit reconstruction with one-shot view prediction, using a completed representation to reduce additional observations of small missing regions \citep{hu2024}. MA-SCVP first obtains one additional NBV and then predicts a remaining view set, using long-tail multiview sampling to improve training \citep{pan2025}. Both approaches are relevant to the complementarity of SPC images and the cost of frequent replanning. For quantitative planetary mapping, an additional distinction is necessary: a plausible completed surface and a surface supported by sufficient measurements have different evidential status. A planner should preserve that distinction when reporting the area meeting an accuracy requirement.

Informative path planning provides a broader connection between map quality and executable motion. \citet{schmid2020} developed an online sampling-based planner that maintains and refines candidate trajectories, with applications including reconstruction of a truncated signed-distance field. Their treatment of information gain and path cost shows how reconstruction feedback can enter a continuing planning process. Applying this principle to SPC requires an observation utility sensitive to photometric and geometric complementarity, as well as feasible spacecraft observation times. It also requires accounting for the delay between acquiring an image and making its reconstruction feedback available.

\subsection{Reinforcement learning for sequential and continuous view planning}

Reinforcement-learning approaches treat observation selection as a sequence whose outcome depends on the accumulated measurements. Scan-RL learns an NBV policy for reconstructing houses and evaluates the reuse of one policy on multiple targets, including houses absent from training \citep{peralta2021}. GenNBV extends learned active reconstruction to a five-dimensional free-space action formulation and combines geometric, semantic, and action representations to improve transfer across datasets \citep{chen2024}. Together, these studies demonstrate that learned observation policies and flexible viewpoint actions have established precedents. Their relevance to SPC lies in the representation of a changing reconstruction state and the optimization of a sequence, while their reported coverage results do not establish planetary terrain accuracy.

Continuous-control methods provide the algorithmic background for learning actions such as camera poses. Deterministic actor--critic methods learn policies over continuous action spaces \citep{lillicrap2015}, while soft actor--critic combines off-policy learning with a stochastic maximum-entropy objective \citep{haarnoja2018}. These approaches illustrate how a policy can produce continuously valued decisions without enumerating a finite action catalogue. They provide a methodological basis for coupling a geometric prediction network to a policy with continuous outputs. In SPC observation planning, this capability must be paired with a reconstruction-specific objective: the desired outcome is a sequence of poses whose images support accurate terrain recovery, with success measured after reconstruction.

For SPC, the sequence objective must capture the complementarity of images. The marginal benefit of an observation depends on which other images already cover the same terrain and on their viewing and illumination conditions. A policy can therefore encounter delayed benefits: a view that is insufficient in isolation may become valuable when a later observation completes a useful combination. This suggests relating learning feedback to progress toward regional reconstruction requirements over successive rounds. Immediate coverage, pose novelty, or image count may provide intermediate signals, but their relation to the final geometric outcome requires evaluation. Otherwise, optimizing the surrogate can encourage repeated acquisition without resolving the uncertainty that actually limits the surface model.

Feedback timing introduces a further distinction from many simulated NBV settings. An acquired image may not yet be available to a ground-based SPC process, and the current model can lag behind execution. The decision state should therefore account for the image history and for which measurements have already contributed to reconstruction, together with the remaining resources and observing opportunities. It is an information state assembled from available evidence; the full true surface is not observed by the planner. This viewpoint helps prevent a simulated policy from relying on knowledge, such as future illumination outcomes or true terrain errors, that would be unavailable during deployment.

Continuous extrinsics must also describe an admissible observation. The position and orientation conventions must be consistent, predicted rotations must represent valid orientations, and the desired pose must be associated with a reachable acquisition time. Spacecraft motion, target rotation, pointing limits, and sunlight jointly restrict these possibilities. Continuous prediction broadens the representation of possible actions but does not make every pose reachable. The planning formulation must therefore distinguish the geometric pose prediction from its admissibility within a mission sequence, including the possibility that further useful acquisition must wait for another opportunity.

Finally, resource constraints need explicit treatment. Constrained policy optimization separates a task reward from auxiliary cost constraints within reinforcement learning \citep{achiam2017}. This provides a useful conceptual precedent for balancing reconstruction progress against resource consumption. However, a bound on an expected cumulative cost does not by itself ensure that a recorder limit is respected at every instant. For spacecraft observation planning, instantaneous occupancy, finite communication windows, and command-delivery timing must be represented at the level where feasibility is checked. The learned policy can optimize scientific decisions within that structure; assigning a negative reward to an overflow event is not equivalent to preventing it.

\subsection{Small-body imaging policies and modeling-aware observation planning}

The closest precedents explicitly connect imaging decisions with small-body mapping. \citet{piccinin2022} formulate acquisition-time selection using deep reinforcement learning, comparing Deep Q-Network and Neural Fitted Q-Iteration policies. Their policy uses a coarse shape model, relative pose, illumination, and acquisition history, with a facet-level mapping index and a limited number of images per storage interval. The evaluation includes bodies with different morphologies, including Eros, Itokawa, Bennu, and 67P. This work already addresses autonomous acquisition, cross-body policy transfer, and data limitations. Two distinctions are relevant to our formulation: the action produced by the modified PointNet describes continuous camera extrinsics, and the intended decision feedback comes from successive SPC reconstructions and their regional quality assessments. Relating observation geometry to locally validated surface error remains central to that feedback.

\citet{wang2025} address observation planning under small-body modeling constraints using a low-resolution model and a genetic algorithm. Their formulation incorporates illumination and shadowing, surface coverage, and operational costs associated with maneuvering and image transmission. This is a direct precedent for planning observations to support photometric reconstruction, and it precludes treating the combination of SPC requirements and resource-aware planning as an unexplored idea. The present study is positioned around a more specific relationship: how deficiencies measured after each reconstruction change the value of candidate image combinations, and how those changes support predictable, verifiable accuracy across different target shapes.

The distinction has consequences for evaluation. A geometry-based mapping score can be a useful planning surrogate, but its improvement must be compared with the geometric improvement actually obtained by SPC. Relevant comparisons therefore include uniform acquisition, fixed SPC observing rules, coverage-oriented selection, and modeling-aware selection under matched opportunities and budgets. For continuous extrinsic prediction, a discretized-action baseline should be given the same reachable region and comparable planning effort; otherwise, gains may reflect a larger admissible action set rather than a better decision rule. Comparing a fixed initial plan with a plan updated from reconstruction diagnostics can then isolate the value of feedback.

There is also a difference between predicting a geometrically appealing pose and learning when that pose is worth acquiring. Similar camera positions can have different utility at different times because the Sun direction, available image set, and recorder occupancy have changed. The PointNet prediction must therefore be interpreted within the sequential planning context. This is the reason for combining geometric pose prediction with reinforcement learning: the former connects the evolving surface to continuous observation geometry, while the latter can optimize the consequences of choosing those observations in a particular order. Establishing that this coupling improves accuracy requires actual SPC reconstructions, including cases where coverage improves but geometric error does not.

\subsection{Mission execution, communication, and storage constraints}

Mission studies illustrate how scientific observation requirements become coupled to time and resources. The Psyche mission design relates mapping activities to its orbit architecture, illumination, stereo imaging needs, and operational margins \citep{polanskey2025}. For Lucy, the Science Encounter Sequence Simulator evaluates commanded observations against surface coverage and resolution requirements and tests sensitivity to uncertainties in encounter geometry and pointing \citep{salmon2025}. These examples support treating the feasibility and robustness of an observation sequence as part of its scientific value. They also motivate checking complete planned sequences, since an attractive individual image may be incompatible with the timing of other required observations.

At the planning-system level, SB-ObsGen combines language-model-based task interpretation with physics-informed tools and hierarchical checks to generate flyby observation schedules \citep{sbobsgen2026}. Its focus on physically and operationally consistent schedules is complementary to the reconstruction-quality problem considered here. Similarly, \citet{castano2022} examine how mission operations can communicate scientific intent to an autonomous spacecraft and recover an understandable account of its decisions. These contributions support automated planning and supervision, while leaving the reconstruction-specific value of an image combination to an appropriate scientific model.

For iterative SPC, resource accounting must follow the lifetime of the data. Acquisition consumes storage immediately, whereas downlink releases capacity according to the communication schedule and the applicable data-retention policy. A ground-based reconstruction can use only data that have arrived; a new ground-generated plan can affect execution only after it has been delivered. Uplink capacity also limits the volume of commands and model updates that can be delivered during a contact. With onboard processing, the data products and exchange requirements differ. Thus, a fixed total-image budget is a useful approximation, but cannot represent every consequence of time-varying uplink, downlink, and recorder occupancy. The present problem treats these quantities as explicit conditions on an adaptive observation campaign, with the computation location left to the mission architecture.

\subsection{Positioning of the proposed PointNet and reinforcement-learning framework}

The proposed framework combines a modified PointNet that predicts continuous camera extrinsics with reinforcement learning for SPC observation planning. Its methodological position follows from three connections. First, the prediction is conditioned on a representation of the reconstructed surface, so observation geometry can respond to the current target morphology. Second, the planning loop is updated by SPC reconstruction and regional quality assessment, connecting subsequent decisions to deficiencies revealed by the preceding image set. Third, the sequence is evaluated within the acquisition, communication, and storage conditions that determine when an image can contribute to a model. These connections organize the earlier research priorities into one observation and reconstruction process.

Each priority also requires its own evidence. Morphological generalization calls for tests on held-out bodies, with consistent reconstruction settings and an explicit account of any target-specific adjustment. Adaptive planning calls for comparisons that remove or delay reconstruction feedback, together with regional checks of the accuracy indicators used in the loop. Resource-limited operation calls for sequence-level accounting of data occupancy and communication, including the consequences of delayed model updates. Ablations of the PointNet modifications and of the continuous pose output can then separate the value of geometric representation from the value of the reinforcement-learning policy.

The central claim to be tested is therefore that this coupling can turn evolving reconstruction needs into useful, executable observation sequences. PointNet predicts the observation extrinsics; the reconstruction and validation components establish what accuracy those observations actually support. This separation makes the framework responsive to the three research objectives while keeping their assessment tied to measurable outcomes. Success is defined by the ability to achieve and document the required surface accuracy across different shapes and constrained campaigns, including an explicit report of regions for which the available observations remain insufficient.
